% Generated draft: requires mathematical and source review.
\documentclass[11pt]{article}
\usepackage[margin=25mm]{geometry}
\usepackage{fontspec}
\setmainfont{Latin Modern Roman}
\usepackage{amsmath,amssymb,mathtools,unicode-math}
\setmathfont{Latin Modern Math}
\usepackage{enumitem,needspace,xurl,hyperref}
\setlength{\parindent}{0pt}
\setlength{\parskip}{5pt}
\setlength{\emergencystretch}{3em}
\newcommand{\N}{\mathbb{N}}
\newcommand{\Z}{\mathbb{Z}}
\newcommand{\Q}{\mathbb{Q}}
\newcommand{\R}{\mathbb{R}}
\newcommand{\C}{\mathbb{C}}
\newcommand{\F}{\mathbb{F}}
\newcommand{\A}{\mathbb{A}}
\newcommand{\PP}{\mathbb{P}}
\newcommand{\E}{\mathbb{E}}
\newcommand{\eps}{\varepsilon}
\newcommand{\dd}{\,\mathrm{d}}
\begin{document}
\section*{UnsolvedMath ID: 30002567}
\textbf{Source problem code:} catalogue entry
\subsection{Definitions and mathematical statement}
The following is the conventional homogeneous real formulation of the supplied conjecture. The source record does not explicitly specify its normalization or boundary conventions; those choices are stated here and require confirmation against the cited report.

Let \(n,d\geq 1\), and let \(h\in\R[x_1,\ldots,x_n]\) be homogeneous of degree \(d\). A polynomial \(h\) is \emph{hyperbolic with respect to} \(e\in\R^n\) if \(h(e)>0\) and, for every \(x\in\R^n\), the polynomial \(t\mapsto h(te-x)\) has only real roots. Here positivity at \(e\) is a normalization convention. Define its closed hyperbolicity cone by
\[
 C_h(e)=\bigl\{x\in\R^n:\text{every root of }t\mapsto h(te-x)\text{ is nonnegative}\bigr\}.
\]
Thus this convention includes the boundary. Apply the same definition to any positive-degree homogeneous polynomial hyperbolic with respect to \(e\). A positive constant polynomial is also allowed as a hyperbolic polynomial of degree zero, with cone \(\R^n\).

A homogeneous polynomial \(p\) of degree \(m\geq1\) has a \emph{definite determinantal representation with respect to \(e\)} if there exist real symmetric matrices \(A_1,\ldots,A_n\in\R^{m\times m}\) such that
\[
 p(x)=\det\!\left(\sum_{i=1}^n x_iA_i\right)
 \quad\text{for every }x\in\R^n,
 \qquad
 \sum_{i=1}^n e_iA_i\succ0.
\]
Here \(\succ0\) means positive definite.

\textbf{Conjecture (supplied formulation of the generalized Lax conjecture).} For every \(n,d\geq1\), every homogeneous \(h\in\R[x_1,\ldots,x_n]\) of degree \(d\), and every \(e\in\R^n\) with respect to which \(h\) is hyperbolic, there exist an integer \(r\geq0\) and a homogeneous polynomial \(q\in\R[x_1,\ldots,x_n]\) of degree \(r\), hyperbolic with respect to the same \(e\), such that
\[
 C_h(e)\subseteq C_q(e)
\]
and \(qh\) has a definite determinantal representation with respect to \(e\). Explicitly, there must exist real symmetric matrices \(A_1,\ldots,A_n\in\R^{(d+r)\times(d+r)}\) satisfying
\[
 q(x)h(x)=\det\!\left(\sum_{i=1}^n x_iA_i\right)
 \quad\text{for every }x\in\R^n,
 \qquad
 \sum_{i=1}^n e_iA_i\succ0.
\]
The inclusion need not be strict. No bound on \(r\) is supplied. Homogeneity, real coefficients, positive normalization, real symmetric matrices, definiteness at \(e\), the closed-cone convention, and allowance of constant multipliers are standard conventions made explicit here; the supplied statement does not define them. Its treatment of constant input polynomials \(h\) is unspecified.
\subsection{Short English statement}
Given a real homogeneous polynomial hyperbolic in a direction \(e\), can one multiply it by another polynomial hyperbolic in the same direction, whose hyperbolicity cone contains the original cone, so that the product is the determinant of a real symmetric linear matrix pencil positive definite at \(e\)? This draft uses explicit standard conventions that require source confirmation. The supplied dated curation labels the problem open; that status has not been independently verified.
\subsection{Sources}
\begin{itemize}
\item Local registry: \texttt{data/unsolvedmath/problems.json}, supplied record with permanent UnsolvedMath ID \(30002567\), title \emph{Definite Determinantal Multiples of Hyperbolic Polynomials}, and exact source problem code \texttt{catalogue entry}. The record attributes the problem to the Oberwolfach Reports Open Problems collection, gives proposed year 2014, and names no individual proposer.
\item Supplied report citation: \emph{Real Algebraic Geometry With A View Toward Systems Control and Free Positivity} (2015), Oberwolfach Reports, pp.\ 977--1045, DOI \texttt{10.4171/owr/2014/17}. The supplied literature-source label is ``Oberwolfach Report 2014/17'': \url{https://doi.org/10.4171/owr/2014/17}. The supplied citation year is retained despite the 2014 report identifier.
\item Supplied literature source, labeled ``Current literature on hyperbolicity cones'': \url{https://doi.org/10.1016/j.indag.2025.06.008}. No authors or article title are supplied.
\item Supplied literature source, labeled ``January 2026 proof claim'': \url{https://arxiv.org/abs/2601.12267}. The supplied assessment describes this as an unrefereed preprint claiming a proof whose claim has not been independently validated.
\item Literature assessment checked at \texttt{2026-08-21}: the supplied curation states that the generalized Lax conjecture remains unresolved in the authoritative literature and retains the status ``open'' despite the January 2026 proof claim. This is supplied dated curation derived from the OWR record, not an independently verified assessment. The record expressly calls for independent verification of time-sensitive status and cited claims.
\end{itemize}
\end{document}
